\section{Choosing an agent framework without changing the causal question}
\label{sec:exec24-framework-census}

An agent framework can help a bank turn a loosely specified business request
into a sequence of inspectable analytical tasks. It can retain a conversation,
route work, call a causal engine and recover an interrupted calculation. The
scientific question remains separate: what intervention is proposed, for whom,
at what information time, under which assumptions, and with what uncertainty?
The useful comparison is therefore between concrete execution contracts and
their ability to preserve that question. A returned sentence, a completed
workflow and a qualified causal recommendation are distinct outcomes
\citep{exec24_framework_contracts}.

Consider a request to increase retention among accounts at risk of leaving.
An assistant can elicit whether the action is a fee waiver, a contact policy
or a credit-limit change; identify the decision date; retrieve a data
dictionary; propose competing graphs; and submit an explicit intervention
query. Every step can be useful before a production product exists. Public
fixtures can test whether the assistant preserves treatment versions,
account histories, unresolved confounding and refusal messages. A later
bank experiment is justified only after these cheaper checks demonstrate
that the proposed chain addresses the relevant scientific problem
\citep{exec24_framework_contracts}.

\subsection{Source profiles and dependencies are part of the comparison}
The inspected LangChain repository contains distinct declarations for
modern LangChain 1.4.3, core 1.6.7 and classic 1.0.8. The modern package
declares a LangGraph range, rather than implementing every execution layer
itself. The separately retained LangGraph source declares 1.2.14.
Comparing the packages as independent causal engines would double-count
the shared orchestration dependency and misstate what either provides
\citep{exec24_framework_langchain,v12_la04_langgraph}.

The other inspected profiles are Semantic Kernel's Python 1.45.0,
DSPy 3.4.0, LlamaIndex core 0.14.25 and a Haystack source declaration of
3.4.0-rc0. A release-candidate label must survive the comparison; it is
not silently promoted to a stable distribution. LlamaIndex's workflow
modules re-export another package. The traced repository resolves to
\texttt{run-llama/llama-agents}, where the selected package declares
\texttt{llama-index-workflows} 2.25.0. This declaration fits core's
specified range, but no solved environment or native compatibility test
has been performed
\citep{exec24_framework_semantic_kernel,exec24_framework_dspy,
exec24_framework_llama_index,exec24_framework_haystack,
exec24_framework_workflows}.

These are authenticated source revisions with commit dates no later than
7 October 2026, inspected from retained original bytes. They are not
claims about the latest release, a bank's installed version, or all
permitted dependency combinations. Root notices for the selected
LangChain, Semantic Kernel, DSPy, LlamaIndex and workflow sources are MIT;
Haystack's selected root notice is Apache License 2.0. AutoGen's earlier
review keeps its AgentChat code notice separate from its documentation
notice. Complete dependency, model-weight, connector and hosted-service
rights remain component-specific inputs
\citep{exec24_framework_langchain,exec24_framework_semantic_kernel,
exec24_framework_dspy,exec24_framework_llama_index,
exec24_framework_haystack,exec24_framework_workflows,v12_la04_autogen}.

\subsection{LangChain and LangGraph: graph recovery is not side-effect recovery}
The modern LangChain factory constructs a state graph and passes a
checkpointer, store and interrupt boundaries into compilation. This is
a concrete route for retaining the business brief, candidate graphs,
assumption decisions and tool results. The factory's configured recursion
limit is a graph-execution control; it is not a token budget, monetary
budget or end-to-end service deadline
\citep{exec24_framework_langchain}.

The retained LangGraph interrupt documentation states that resumption
restarts the node and re-executes its logic. A checkpointer is required
for the documented interrupt route. These properties support recoverable
analysis, but preceding external actions need their own replay contract.
If a node creates a paid external job before the interrupt, restored graph
state does not by itself prove that the job can safely be created again
\citep{v12_la04_langgraph,exec24_framework_contracts}.

The inspected modern retry middleware makes a useful further distinction.
It counts the initial attempt plus the configured retries, propagates
graph control-flow exceptions, and treats nonmatching exceptions separately.
After exhaustion, the tool middleware can return a tool message explicitly
marked as an error; the model middleware can return an assistant message
containing failure information. Both also permit an error-raising route.
A wrapper that tests only whether a message exists would erase the
difference between recovered output and reported failure
\citep{exec24_framework_langchain}.

An authored banking example specifies the resulting public test. Suppose
an external service accepts a 500 USD action, but its acknowledgement is
lost. Three submissions without a shared deduplication key can represent
1,500 USD of actions. Under an explicitly stipulated service contract
that deduplicates the same key, they represent 500 USD. The 1,000 USD
difference is not an observed framework defect; it is the consequence
of different external contracts. The proposed adapter must preserve the
logical action key, attempt identity, acknowledgement state and recovery
query. Retrying an unidentified causal request also does not repair its
scientific assumptions \citep{exec24_framework_contracts}.

LangChain and LangGraph are retained as reference orchestration choices,
with conditional integration through typed causal-engine tools. Direct
adoption as the scientific product is deferred. Porting a small source
mechanism is an alternative only when its dependency and notice obligations
are understood and its semantics can be maintained. A custom controller
may be simpler for a narrowly bounded workflow; the comparison should
measure that alternative on the same recovery and refusal fixtures
\citep{exec24_framework_contracts}.

\subsection{AutoGen: conversational state and scientific state are different}
The retained AutoGen AgentChat 0.7.5 source says that its assistant retains
state across calls and should receive new messages rather than the whole
conversation again. It warns against concurrent calls on the same agent
instance. The documented tool behavior can involve concurrent tool calls
and an additional model call for reflection. Its save/load methods retain
model context; they do not establish transactional recovery of arbitrary
external side effects
\citep{v12_la04_autogen}.

For the retention example, separate conversational roles can propose
assumptions, challenge them and explain results. Agreement among those
roles is not independent scientific evidence if they share the same
model, data, prompt history or mistaken business premise. Their graph
proposals should enter a comparison set with provenance, not become
accepted orientations because a conversation reached consensus. Scientific
state must include the graph object, query, assumptions, identification
result and uncertainty object independently of the natural-language history
\citep{exec24_framework_contracts}.

The pinned README declares maintenance mode and recommends a successor.
That declaration informs a maintenance decision; the successor's advertised
properties are not established here by repeating the README. Reference
use and an existing-installation adapter remain possible, while a new
deployment decision needs an independently reviewed migration target and
matched fixtures. The proposal therefore defers unconditional adoption,
retains a conditional integration route, and records maintenance and
component-specific rights as explicit revisit conditions
\citep{v12_la04_autogen,exec24_framework_contracts}.

\subsection{Semantic Kernel: function filters do not verify scientific claims}
The inspected Python function-choice behavior exposes included/excluded
plugins and functions, and a maximum automatic-invocation setting. Its
kernel checks required and unexpected argument names, constructs an
invocation-filter stack, and can return a termination context. These are
useful controls for exposing only the intended causal tools at a given
stage. They do not verify the truth of a proposed exclusion restriction
or the adequacy of an adjustment set
\citep{exec24_framework_semantic_kernel}.

The selected automatic-invocation handler catches ordinary function
exceptions and places error text in a function result. A separate guard
is therefore needed before a later assistant message is accepted as a
qualified answer. The proposed protocol distinguishes an invalid request,
an unavailable service, a scientific refusal, a partial numerical result
and a successful calculation with stated assumptions. Argument-schema
validity is one condition among these, not their replacement
\citep{exec24_framework_semantic_kernel,exec24_framework_contracts}.

This review covers the selected Python profile. It does not transfer its
behavior to .NET, Java or every provider connector. Conditional integration
would expose a narrow tool set, retain scientific output objects and
test filters, failure results and termination using public fixtures.
Reference use is retained; broad cross-language parity and direct adoption
as the causal library are deferred. Porting requires separately inspecting
the actual component and maintaining its conversion contract
\citep{exec24_framework_contracts}.

\subsection{DSPy: optimizing a program requires a defensible evaluator}
The inspected ReAct implementation builds a trajectory, calls selected
tools and records execution exceptions as observations. It may leave the
tool loop after an invalid selection or context-window problem and then
run an extraction step. This is a useful mechanism for explaining partial
work, but the extraction step is not proof that the requested causal
calculation succeeded. The result must retain whether identification,
estimation and uncertainty checks actually completed
\citep{exec24_framework_dspy}.

The evaluation implementation supplies a particularly relevant positive
contract: after a returned evaluation run, missing per-example results
receive the configured failure score, and the aggregate denominator is
the full supplied development set. That guards against silently reporting
only successful examples in that code path. If the executor raises at
its error threshold, however, no completed aggregate should be invented.
The task definition, scorer, allowed failure score and failure threshold
must accompany every reported benchmark
\citep{exec24_framework_dspy}.

In an authored example, 1,000 requests produce 800 returned answers and
600 scientifically qualified answers. The qualified fraction among
returns is $600/800=3/4$, while the fraction among all requests is
$600/1000=3/5$. These answer different questions. Neither fraction alone
measures causal calibration, utility or the severity of unqualified
answers. A public experiment should report the complete request ledger,
refusals, errors, qualified answers and the prespecified loss assigned to
each outcome \citep{exec24_framework_contracts}.

The inspected parallel executor also deserves careful source-level
qualification. It can resubmit stragglers and stores completed outcomes
by input index, while a resubmission set is keyed by Future objects.
The nearby comment describing one resubmission per item is therefore not
by itself a proved per-item invariant: distinct Future identities can
refer to the same input. No native liveness or duplicate-work failure is
asserted here. A public test must trace the input key and every submission,
including work that continues after a useful result has been collected
\citep{exec24_framework_dspy,exec24_framework_contracts}.

DSPy is retained as a reference for program construction and evaluation,
with conditional integration around a fixed causal-task scorer. The
scorer must distinguish graph prediction, identification, numerical
estimation, uncertainty and business loss. Optimizing on the same examples
later used to claim performance would change the evaluation question.
Adoption as a causal oracle is rejected; bounded program optimization and
selective source reuse remain testable alternatives
\citep{exec24_framework_contracts}.

\subsection{LlamaIndex: stopping, structured output and runtime delegation}
The selected LlamaIndex base agent keeps iteration counts in context.
At its limit, the inspected path either raises or performs a final model
generation, depending on the chosen early-stopping method. A generated
summary after exhausting the budget is useful communication, but its
scientific status must remain partial or unresolved when necessary.
The selected function-agent code also distinguishes ordinary tool output
from a return-direct route and from handoff behavior
\citep{exec24_framework_llama_index}.

The separate workflows package makes the time contract especially clear.
Its selected constructor documents timeout in known-alive seconds accrued
through processed journal ticks; downtime during resumption and some
pre-crash time do not count. This is a different clock from an external
bank service deadline. In an authored example, two 20-second active
periods separated by a one-hour outage total 40 active seconds but
3,640 elapsed seconds. A caller needing a 60-second wall deadline cannot
infer it from the active-time limit
\citep{exec24_framework_workflows,exec24_framework_contracts}.

The inspected cancellation method signals the runtime and waits for the
result, while catching timeout, cancellation and other exceptions. Its
return therefore does not independently prove that every remote side
effect has stopped. The retry-policy source separately defines retry
predicates, delays and stopping conditions. These concrete distinctions
support an adapter design with external job reconciliation, explicit
deadline clocks and durable attempt identifiers. They do not establish
that the complete distributed runtime has been qualified
\citep{exec24_framework_workflows,exec24_framework_contracts}.

Conditional integration is attractive when document retrieval and a
workflow context are needed together, provided retrieved statements stay
traceable to inspected source material. Reference use is retained and
direct adoption as the scientific layer is deferred. The actual delegated
runtime, serializer, persistence backend and model connector need their
own fixture results. A source-level range match is only the start of
that engineering comparison \citep{exec24_framework_contracts}.

\subsection{Haystack: preserve explicit stop reasons and hook effects}
The selected Haystack agent supports state mappings, ordered hooks,
tool-concurrency control and a maximum number of agent steps. Its
documented step consists of a generator call and the requested tool work.
The inspected limit path records a maximum-steps exit reason and returns
the current public state. A tool-based exit checks error information;
the selected logic does not treat an errored exit-condition tool as a
successful exit. These are valuable details to preserve when comparing
returned results
\citep{exec24_framework_haystack}.

Hooks are scientifically consequential. A before-tool hook can alter
the pending request; an after-tool hook can transform returned messages;
an after-run hook can alter the final state. These facilities can reduce
context size or enforce a protocol, but a summary must not silently
discard an unresolved edge, confidence set or refusal. The selected
general function-based Tool class also lacks a formal output-name schema
by default, whereas specialized components can provide one. The proposed
adapter must therefore define and validate the causal result itself
\citep{exec24_framework_haystack,exec24_framework_contracts}.

For an authored example, the uncertainty sets $[-1,1]$ and $[-2,2]$
both have midpoint zero. A hook that keeps only the midpoint makes them
indistinguishable and cannot reconstruct either set from that scalar.
Whether a bank action is acceptable can depend on the endpoints even
when the midpoint is identical. Typed results should preserve the
original set, units, assumptions and numerical diagnostics; prose can
refer to that object without replacing it
\citep{exec24_framework_contracts}.

Haystack remains a reference and conditional pipeline choice. The
release-candidate source identity, tool schemas, hooks, state mappings
and runtime dependencies require matched public tests. Blanket adoption
as a causal identification engine is rejected. Selective reuse or a
smaller custom pipeline can be compared on the same preservation and
failure criteria \citep{exec24_framework_contracts}.

\subsection{A common public qualification design}
The comparisons suggest a shared test design rather than six incompatible
demonstrations. Start with the same business brief, data dictionary,
candidate graph set, formal query and causal-engine oracle. Change one
framework or backend factor at a time where attribution is intended;
retain the full factorial design when interactions matter. Test both
successful and deliberately unsupported cases: hidden-confounding
ambiguity, off-support interventions, unknown tool outcomes, malformed
schemas, interrupted state, duplicated acknowledgements and expired
deadlines. A persuasive narrative is not the scoring oracle
\citep{exec24_framework_contracts}.

Concurrency requires a testable budget invariant. If two workers each
read an unreserved balance of 100 USD and each accepts a 70 USD task,
both local checks pass while the combined commitment is 140 USD.
An atomic reservation rule can prevent that particular overspend;
a language-model instruction to stay within budget cannot establish
atomicity. The fixture must specify the transaction and recovery model,
then test reservations, release and repeated requests under that model
\citep{exec24_framework_contracts}.

The common adapter contract should carry immutable task and graph
identities, intervention and target, source citations, assumptions,
information timestamps, scientific status, numerical output, uncertainty,
logical action keys, attempts, deadlines and complete costs. It should
preserve uncertainty and refuse unsupported conversions. A CLI and an MCP
interface can transport this same contract; transport success does not
change the scientific verdict. Integrations with other repositories should
use their declared interfaces and public fixtures, with changes confined
to the proposed library's adapter until separate work is authorized
\citep{exec24_framework_contracts}.

Scientific ownership belongs to the target and assumptions; statistical
ownership to evaluation and dependence; adapter ownership to types,
refusals and replay; platform ownership to exact environments, rights,
budgets and recovery. Under hypothetical expenditure of 24,000 USD,
1,000 submitted, 800 returned and 600 qualified requests cost 24, 30 and
40 USD respectively. Actual provider prices, latency, native reliability
and bank value remain unmeasured
\citep{exec24_framework_contracts}.

The bounded reference choices are now concrete enough to guide a public
prototype comparison. They do not close the hosted/local backend survey,
the remaining framework and connector delegates, or the frontier literature
census. No framework removes unobserved confounding or recovers discarded
information without additional assumptions. Package execution and independent validation remain to be performed
\citep{exec24_framework_contracts}.
